% Generated draft: requires mathematical and source review.
\documentclass[11pt]{article}
\usepackage[margin=25mm]{geometry}
\usepackage{fontspec}
\setmainfont{Latin Modern Roman}
\usepackage{amsmath,amssymb,mathtools,unicode-math}
\setmathfont{Latin Modern Math}
\usepackage{enumitem,needspace,xurl,hyperref}
\setlength{\parindent}{0pt}
\setlength{\parskip}{5pt}
\setlength{\emergencystretch}{3em}
\newcommand{\N}{\mathbb{N}}
\newcommand{\Z}{\mathbb{Z}}
\newcommand{\Q}{\mathbb{Q}}
\newcommand{\R}{\mathbb{R}}
\newcommand{\C}{\mathbb{C}}
\newcommand{\F}{\mathbb{F}}
\newcommand{\A}{\mathbb{A}}
\newcommand{\PP}{\mathbb{P}}
\newcommand{\E}{\mathbb{E}}
\newcommand{\eps}{\varepsilon}
\newcommand{\dd}{\,\mathrm{d}}
\begin{document}
\section*{UnsolvedMath ID: 30002664}
\textbf{Source problem code:} catalogue entry
\subsection{Definitions and mathematical statement}
Let \(\mathcal R\) denote the intended class of planar regions; the supplied record does not define this class. Write \(\operatorname{Area}(R)\) for planar area, \(\operatorname{int}(R)\) for the topological interior of \(R\subseteq\R^2\), and \(\Z^2=\{(a,b):a,b\in\Z\}\) for the standard integer lattice. For any set \(S\), let \(\#S\) denote its cardinality, which may be infinite. Boundary lattice points are excluded from the count.

Under the usual interpretation of congruence as an arbitrary Euclidean isometry, a congruent copy of \(R\) has the form \(QR+t\), where \(Q\in O(2)=\{Q\in\R^{2\times2}:Q^{\mathsf T}Q=I_2\}\) and \(t\in\R^2\). Thus reflections as well as rotations and translations are included.

The supplied question asks whether
\[
\forall n\in\N\quad \exists R\in\mathcal R\quad
\left[\operatorname{Area}(R)=n\ \land\
\forall Q\in O(2)\ \forall t\in\R^2,\quad
\#\bigl(\operatorname{int}(QR+t)\cap\Z^2\bigr)\ne n\right].
\]
Here each region may depend on \(n\). The source does not specify whether \(0\in\N\), or whether regions must be measurable, bounded, connected, open, closed, or satisfy boundary regularity conditions. Consequently, \(\mathcal R\) and the meaning of area remain unresolved qualifications of this statement. If area is interpreted as Lebesgue measure, measurability must be specified. The standard lattice and the full Euclidean congruence group are conventional interpretations rather than explicitly supplied definitions.
\subsection{Short English statement}
For every natural number \(n\), can one find a planar region of area \(n\) such that every congruent placement has an interior lattice-point count different from \(n\)? Points on the boundary do not count. The admissible meaning of planar region and the convention concerning \(n=0\) require clarification.
\subsection{Sources}
\begin{itemize}
\item Local UnsolvedMath registry, supplied record with permanent ID \(30002664\), title \emph{Planar Regions Avoiding Prescribed Interior Lattice Point Counts}, exact source problem code \texttt{catalogue entry}. The record labels the statement exact and describes it as retained byte-for-byte; these are supplied registry annotations, not independent source verification.
\item Supplied source citation: \emph{Discrete Geometry (2015)}, Oberwolfach Reports, pp.\ 2235--2296, DOI \texttt{10.4171/owr/2014/40}. \url{https://doi.org/10.4171/owr/2014/40}. The record gives proposed year 2014 and supplies no named proposer. Its background describes Oberwolfach Reports as documenting mathematical workshops and their open-problem discussions.
\item Supplied literature source: \emph{Oberwolfach report on the Steinhaus lattice-point problem}. \url{https://ems.press/content/serial-article-files/46530}.
\item Supplied literature source: \emph{Maehara, lattice points on polyhedra}. \url{https://doi.org/10.1007/s00454-017-9869-8}.
\item Literature assessment dated exactly \texttt{2026-08-21}: the supplied curation labels the problem open and reports positive results for compact, strictly convex, algebraic, and polygonal regions, while describing the arbitrary-region question as unresolved. This is supplied dated curation derived from the OWR material, not an independently verified assessment. Its status and cited claims require independent verification.
\end{itemize}
\end{document}
